\hypertarget{native-282-nexus6}{}\label{native-282-nexus6}
\Article{6}
Before investigating the consequences in general, we have thought it useful to work out some special cases, so that the manner of argument may be understood more easily. Let us therefore first return to the form $xx+yy$, for which the number of representations of an indeterminate number has the mean value $=\pi$. The number of actual representations of a given number is readily determined from the general principles established in the Arithmetical Investigations. Let $fA$ denote the number of representations of the number $A$; this will be $=4$ if $A=1$ or $2$ or a power of two; $=8$ if $A$ is a prime number of the form $4n+1$, or the product
\SourceLeaf{283}
of such a prime and a power of two; $=0$ if $A$ is a prime number of the form $4n+3$, or is divisible by such a prime but not by its square; finally, \emph{in general},
\[
\begin{aligned}
\text{either}\quad&=4(\alpha+1)(\beta+1)(\gamma+1)\ldots\\
\text{or}\quad&=0
\end{aligned}
\]
Here the number $A$ is put in the form $2^\mu Sa^\alpha b^\beta c^\gamma\ldots$, where $a,b,c$, etc., denote different prime numbers of the form $4n+1$, and $S$ the product of the primes of the form $4n+3$ that occur once or more among the factors of the number $A$, if there are any. The first or second value applies according as the number $S$ is a square or is not a square. It is therefore evident that $fA$ depends solely on the way in which the prime numbers $3,5,7,11,13$, etc., occur among the factors of the number $A$, so that in general we must put
\[
fA=4(3)\pdot(5)\pdot(7)\pdot(11)\pdot(13)\dotsfour
\]
\begin{samepage}
if we suppose the values of the symbols $(3),(5),(7)$, etc., to be chosen so that, with $p$ denoting a prime number, we have

\begin{description}
\item[first] $(p)=1$ if $p$ does not divide $A$
\item[second] $(p)=\alpha+1$ if $p$ has the form $4n+1$, and $p^\alpha$ is the highest power dividing $A$
\item[third] $(p)=0$ if $p$ has the form $4n+3$, and the exponent of the highest power of $p$ dividing $A$ is odd; finally,
\item[fourth] $(p)=1$ if $p$ has the form $4n+3$, and the exponent of the highest power of $p$ dividing $A$ is even.
\end{description}
\end{samepage}
Evidently the first case is included in the second and fourth.

The terms of the sequence $f1,f2,f3,f4$, etc., thus proceed very irregularly, although the more terms are taken, the more accurately the mean value $=\pi$ must emerge from them. We shall denote the sum $f1+f2+f3+\ldots+fA$ by $FA$.

\Article{7}
Let us now put, in general, $fm+f3m=f'm$; it is easily seen that
\[
f'A=4(5)\pdot(7)\pdot(11)\pdot(13)\ldots
\]
\SourceLeaf{284}
that is, $f'A$ will be independent of the relation of $A$ to the divisor $3$, so that the irregularity of the sequence $f'1,f'2,f'3,f'4,f'5,f'6$, etc., will both begin later and be much smaller. Further, if we put
\[
f'1+f'2+f'3+f'4+\text{etc.}+f'm=F'm
\]
we shall have
\[
\begin{aligned}
F'3A&=F3A+f3+f6+f9+\ldots+f9A\\
&=F3A+FA
\end{aligned}
\]
From this it follows readily that, as $A$ increases to infinity, we must put
\[
F'3A=4\pi A
\]
or that the mean value of the terms of the sequence $f'1,f'2,f'3,f'4$, etc., is
\[
=\tfrac43\pi
\]

In the same way, on putting in general $-f'm+f'5m=f''m$, we obtain
\[
f''A=4(7)(11)(13)\ldots
\]
that is, the new sequence $f''1,f''2$, etc., loses the fluctuations that depend on the relation to the number $5$. On setting the sum
\[
f''1+f''2+f''3+\ldots+f''m=F''m
\]
we obtain
\[
F''5m=-F'm+F'5m
\]
from which it follows that, as $m$ increases to infinity, we must put
\[
F''5m=\tfrac43\pi\pdot4m
\]
or that the mean value of the terms of the sequence is $=\tfrac43\pdot\tfrac45\pi$.

If we proceed further in the same way, forming new sequences as we successively remove the factors $(7),(11),(13),(17)$, etc., these will approach constancy ever more closely, and their mean values will successively acquire the new factors $\tfrac87,\tfrac{12}{11},\tfrac{12}{13},\tfrac{16}{17},\tfrac{20}{19}$, etc. The denominators will be the prime numbers in their natural order, and the numerators will be greater or smaller by one according as those primes have the form $4n-1$ or ${}^{.}4n+1$. Hence, since, on continuing this process to infinity, the constant value
\SourceLeaf{285}
$4$ must approach the mean value ever more closely, we have
\[
4=\pi\pdot\tfrac43\pdot\tfrac45\pdot\tfrac87\pdot\tfrac{12}{11}\pdot\tfrac{12}{13}\ldots\text{ to infinity}
\]
or
\[
\pi=4\pdot\frac3{3+1}\pdot\frac5{5-1}\pdot\frac7{7+1}\pdot\frac{11}{11+1}\pdot\frac{13}{13-1}\ldots
\]
If the individual fractions are expanded into infinite series,
\[
\begin{aligned}
\frac3{3+1}&=1-\frac13+\frac19-\frac1{27}+\ldots\\
\frac5{5-1}&=1+\frac15+\frac1{25}+\frac1{125}+\ldots\\
\frac7{7+1}&=1-\frac17+\frac1{49}-\frac1{343}+\ldots\\
&\text{etc.}
\end{aligned}
\]
the product readily expands into
\[
1-\tfrac13+\tfrac15-\tfrac17+\tfrac19-\tfrac1{11}+\tfrac1{13}\dotsfour
\]
and the sum of this series is generally known to be $=\tfrac14\pi$. In fact, the illustrious \textsc{Euler} had long ago derived from this, in the reverse direction, the equality between $\tfrac14\pi$ and the infinite product $\tfrac43\pdot\tfrac45\pdot\tfrac87\ldotp\ldotp$ (\emph{Introd. in analys. inf.}, Vol. I, Chap. xv, Article~285).

\clearpage
\Article{8}
Let us consider, secondly, the form $xx+2yy$, for which the number of representations of an indeterminate number will have the mean value $=\frac\pi{\surd2}$. If $fA$ denotes the number of representations of a given number $A$ by this form, this will be $=2$ for $A=1$ or $A=2$, or whenever $A$ is a power of two; further, $fA=4$ whenever $A$ is one of the primes of which $-2$ is a quadratic residue, that is, those of the forms $8k+1,8k+3$, for example $A=3,11,17,19,41,43$, etc.; finally, $fA=0$ whenever $A$ is a prime of which $-2$ is a quadratic nonresidue, for example one of the sequence $5,7,13,23,29,31$, etc., that is, of either the form $8k+5$ or the form $8k+7$. In general, however, we must put
\[
\begin{aligned}
\text{either}\quad fA&=2(\alpha+1)(\beta+1)(\gamma+1)\ldots\\
\text{or}\quad fA&=0
\end{aligned}
\]
Here the number $A$ is put in the form $2^\mu Sa^\alpha b^\beta c^\gamma\ldots$, where $a,b,c$, etc., denote different prime numbers of the forms $8k+1,8k+3$, while $S$ is the product of the
\SourceLeaf{286}
remaining numbers (of the forms $8k+5,8k+7$), if any occur among the factors of the number $A$. The first or second value applies according as $S$ is a square or is not a square. From this, by reasoning entirely similar to that in the preceding article, we shall pass from the sequence $f1,f2,f3,f4,f5$, etc., namely $2,2,4,2,0,2$, etc., successively to others that remain \emph{constant} for increasingly long stretches. Their \emph{mean} values will be, successively, $\frac\pi{\surd2}$, $\frac\pi{\surd2}\pdot\frac23$, $\frac\pi{\surd2}\pdot\frac23\pdot\frac65$, $\frac\pi{\surd2}\pdot\frac23\pdot\frac65\pdot\frac87$, etc.; we shall proceed in such a way as to arrive at the equation
\[
2=\frac\pi{\surd2}\pdot\frac23\pdot\frac65\pdot\frac87\pdot\frac{10}{11}\pdot\frac{14}{13}\pdot\frac{16}{17}\pdot\frac{18}{19}\ldots
\]
where the denominators form the natural sequence of prime numbers, while the numerators are smaller by one whenever the denominators have the form $8k+1$ or $8k+3$, and larger by one whenever the denominators have the form $8k+5$ or $8k+7$.
\BookRule{.30\textwidth}
